\section{Preliminaries}\label{sec:preliminaries}

\subsection{Embodied Agent Environments}

We consider an embodied agent that interacts with its environment through a minimal interface. At each timestep $t$, the agent receives a frame observation $o_t$ (a rendered image of the current environment state) together with a text map $m_t$ that describes the visible tiles and nearby walkable positions in ASCII form, and selects an action $a_t$ from a fixed set of button inputs $\mathcal{A}$. The text map is derived from game state that a human player can read off the screen, and compensates for the limited spatial reasoning of current vision-language models; it contains no walkthrough, no objective list, and no pathfinding. The environment is partially observable since the agent cannot access internal state such as NPC intent or battle mechanics beyond what the frame and map expose.

\subsection{Agentic Harnesses}

An \emph{agentic harness} $\mathcal{H}$ is the scaffolding layer between a foundation model $M$ and the environment. Following the decomposition from~\citet{karten2026pokeagent}, a harness mediates agent behavior through four components:
\begin{itemize}
    \item \textbf{System prompt} $p$: the instructions and strategic guidance provided to the model at each reasoning step.
    \item \textbf{Sub-agents} $\mathcal{G}$: specialized modules that can be invoked by the orchestrator for specific tasks (e.g., battle strategy, puzzle solving, self-reflection).
    \item \textbf{Skills} $\mathcal{K}$: reusable routines available to the model, spanning both text-level behaviors (heuristics cited in reasoning) and executable programs (pathfinders, tool wrappers). Pre-built primitives such as \texttt{press\_buttons} and \texttt{get\_game\_state} are skills the harness ships with; new skills can also be authored during play.
    \item \textbf{Memory} $\mathcal{M}$: a persistent knowledge store that accumulates facts, strategies, and observations across the agent's trajectory.
\end{itemize}
In addition to these refined components, the harness exposes a fixed set of \emph{meta-tools} (\texttt{define\_agent}, \texttt{run\_code}, \texttt{process\_memory}, and similar primitives) through which the agent edits $p, \mathcal{G}, \mathcal{K}, \mathcal{M}$ in place.

A \emph{minimalist harness} $\mathcal{H}_{\min}$ provides only the environment interface ($o_t$, $m_t$, $a_t \in \mathcal{A}$) with a generic system prompt and no sub-agents, memory, or authored skills. A \emph{hand-engineered harness} $\mathcal{H}_\mathrm{expert}$ populates all components through manual engineering. A meta-harness gives the model meta-tools (\texttt{define\_agent}, \texttt{run\_code}, etc.) to construct its own sub-agents, skills, and memory entries during play; this was the operating point of our later GPP runs, where the model built its own pathfinders, battle strategists, and reusable scripts without being asked to. \texttt{Continual Harness} starts from a minimal harness and adds an automated Refiner that rewrites $p, \mathcal{G}, \mathcal{K}, \mathcal{M}$ in place from trajectory analysis. We write $\mathcal{H}_\mathrm{CH}$ for the running harness state during a \texttt{Continual Harness} run, evolving with every refinement cycle.